% GENERATED FILE. Edit canonical lessons, then run npm run generate:books.
% canonical-source-hash: fnv1a64:b7e82f7ec30bea49
% canonical-lessons: KA-C225-bevu, KA-C225-mallige, KA-C225-suryakanti, KA-C225-kabbu, KA-C225-pataki

\chapter{Neem, Jasmine, Sunflower, Sugarcane, Firecracker}
\label{ch:ka-neem-jasmine-sunflower-sugarcane-firecracker}
\hlchaptermodality{225}

\begin{chapteropening}
\textbf{By the end of this chapter:} \emph{I can say neem, jasmine, a sunflower, sugarcane and a firecracker.}

Complete the last lesson of chapter 225: I can say neem, jasmine, a sunflower, sugarcane and a firecracker.
\end{chapteropening}

\section[\texorpdfstring{\kn{ಬೇವು} (bēvu)}{ಬೇವು (bēvu)}]{\kn{ಬೇವು} (bēvu) — neem}

\label{lesson:KA-C225-bevu}



\begin{quote}
Before the new one: say the Kannada for nature, then the Kannada for an animal.
\end{quote}

\subsection*{You'll want to know: \kn{ಬೇವು}}
\textbf{\kn{ಬೇವು}} — \emph{bēvu} — \textquotedblleft{}neem\textquotedblright{}.

\begin{grammarlens}[title={What you've built}]
One more word for this chapter.
\end{grammarlens}

\subsection*{Your turn}
\begin{itemize}
\raggedright
  \item \emph{Say it:} \emph{bēvu}
  \item \emph{Say it:} \emph{bēvu}, once more
  \item \emph{Say it:} say \emph{bēvu}
  \item \emph{Recall:} say the Kannada for nature, then the Kannada for an animal, then say \emph{bēvu} again
\end{itemize}

\begin{tcolorbox}[breakable,colback=teal!4,colframe=teal!35!black,title={Before you move on}]
What does \kn{ಬೇವು} mean? (\textquotedblleft{}Neem\textquotedblright{}.) Say it once more.
\end{tcolorbox}
\section[\texorpdfstring{\kn{ಮಲ್ಲಿಗೆ} (mallige)}{ಮಲ್ಲಿಗೆ (mallige)}]{\kn{ಮಲ್ಲಿಗೆ} (mallige) — jasmine}

\label{lesson:KA-C225-mallige}



\begin{quote}
Before the new one: say the Kannada for an animal, then the Kannada for neem.
\end{quote}

\subsection*{You'll want to know: \kn{ಮಲ್ಲಿಗೆ}}
\textbf{\kn{ಮಲ್ಲಿಗೆ}} — \emph{mallige} — \textquotedblleft{}jasmine\textquotedblright{}.

\begin{grammarlens}[title={What you've built}]
One more word for this chapter.
\end{grammarlens}

\subsection*{Your turn}
\begin{itemize}
\raggedright
  \item \emph{Say it:} \emph{mallige}
  \item \emph{Say it:} \emph{mallige}, once more
  \item \emph{Say it:} say \emph{mallige}
  \item \emph{Recall:} say the Kannada for an animal, then the Kannada for neem, then say \emph{mallige} again
\end{itemize}

\begin{tcolorbox}[breakable,colback=teal!4,colframe=teal!35!black,title={Before you move on}]
What does \kn{ಮಲ್ಲಿಗೆ} mean? (\textquotedblleft{}Jasmine\textquotedblright{}.) Say it once more.
\end{tcolorbox}
\section[\texorpdfstring{\kn{ಸೂರ್ಯಕಾಂತಿ} (sūryakānti)}{ಸೂರ್ಯಕಾಂತಿ (sūryakānti)}]{\kn{ಸೂರ್ಯಕಾಂತಿ} (sūryakānti) — a sunflower}

\label{lesson:KA-C225-suryakanti}



\begin{quote}
Before the new one: say the Kannada for neem, then the Kannada for jasmine.
\end{quote}

\subsection*{You'll want to know: \kn{ಸೂರ್ಯಕಾಂತಿ}}
\textbf{\kn{ಸೂರ್ಯಕಾಂತಿ}} — \emph{sūryakānti} — \textquotedblleft{}a sunflower\textquotedblright{}.

\begin{grammarlens}[title={What you've built}]
One more word for this chapter.
\end{grammarlens}

\subsection*{Your turn}
\begin{itemize}
\raggedright
  \item \emph{Say it:} \emph{sūryakānti}
  \item \emph{Say it:} \emph{sūryakānti}, once more
  \item \emph{Say it:} say \emph{sūryakānti}
  \item \emph{Recall:} say the Kannada for neem, then the Kannada for jasmine, then say \emph{sūryakānti} again
\end{itemize}

\begin{tcolorbox}[breakable,colback=teal!4,colframe=teal!35!black,title={Before you move on}]
What does \kn{ಸೂರ್ಯಕಾಂತಿ} mean? (\textquotedblleft{}A sunflower\textquotedblright{}.) Say it once more.
\end{tcolorbox}
\section[\texorpdfstring{\kn{ಕಬ್ಬು} (kabbu)}{ಕಬ್ಬು (kabbu)}]{\kn{ಕಬ್ಬು} (kabbu) — sugarcane}

\label{lesson:KA-C225-kabbu}



\begin{quote}
Before the new one: say the Kannada for jasmine, then the Kannada for a sunflower.
\end{quote}

\subsection*{You'll want to know: \kn{ಕಬ್ಬು}}
\textbf{\kn{ಕಬ್ಬು}} — \emph{kabbu} — \textquotedblleft{}sugarcane\textquotedblright{}.

\begin{grammarlens}[title={What you've built}]
One more word for this chapter.
\end{grammarlens}

\subsection*{Your turn}
\begin{itemize}
\raggedright
  \item \emph{Say it:} \emph{kabbu}
  \item \emph{Say it:} \emph{kabbu}, once more
  \item \emph{Say it:} say \emph{kabbu}
  \item \emph{Recall:} say the Kannada for jasmine, then the Kannada for a sunflower, then say \emph{kabbu} again
\end{itemize}

\begin{tcolorbox}[breakable,colback=teal!4,colframe=teal!35!black,title={Before you move on}]
What does \kn{ಕಬ್ಬು} mean? (\textquotedblleft{}Sugarcane\textquotedblright{}.) Say it once more.
\end{tcolorbox}
\section[\texorpdfstring{\kn{ಪಟಾಕಿ} (paṭāki)}{ಪಟಾಕಿ (paṭāki)}]{\kn{ಪಟಾಕಿ} (paṭāki) — a firecracker}

\label{lesson:KA-C225-pataki}



\begin{quote}
Before the new one: say the Kannada for a sunflower, then the Kannada for sugarcane.
\end{quote}

\subsection*{You'll want to know: \kn{ಪಟಾಕಿ}}
\textbf{\kn{ಪಟಾಕಿ}} — \emph{paṭāki} — \textquotedblleft{}a firecracker\textquotedblright{}.

\begin{grammarlens}[title={What you've built}]
One more word for this chapter.
\end{grammarlens}

\subsection*{Your turn}
\begin{itemize}
\raggedright
  \item \emph{Say it:} \emph{paṭāki}
  \item \emph{Say it:} \emph{paṭāki}, once more
  \item \emph{Say it:} say \emph{paṭāki}
  \item \emph{Recall:} say the Kannada for a sunflower, then the Kannada for sugarcane, then say \emph{paṭāki} again
\end{itemize}

\begin{tcolorbox}[breakable,colback=teal!4,colframe=teal!35!black,title={Before you move on}]
What does \kn{ಪಟಾಕಿ} mean? (\textquotedblleft{}A firecracker\textquotedblright{}.) Say it once more.
\end{tcolorbox}
